\chapter{The casualty}\label{ch:casualty}

Everything in this chapter is decision support for trained medics. It organises what was reported, keeps
timers and suggests an order. It never diagnoses or prescribes. Every threshold below is a placeholder to
validate with a medical mentor; all of them sit at the top of their module.

\section{MIST card and priority}

\code{casualty.py} keeps one card per casualty in the format medics use for handover:
\textbf{M}echanism (blast, gunshot, fall, vehicle), \textbf{I}njuries in the speaker's own words,
\textbf{S}igns (conscious, breathing, massive bleeding, AVPU), \textbf{T}reatments with times. Follow-up
reports update the same card. ``He is now unconscious'' changes the signs without the word ``casualty''.

Priority is rule-based. A casualty is \emph{urgent} with massive bleeding, a tourniquet, unconsciousness
or a breathing problem; otherwise \emph{priority} or \emph{routine} from what was said. The evacuation order
is: priority, then who gets worse first, then the oldest tourniquet, then distance to the CCP.

\section{Condition and ``worse in'' windows}

\code{condition.py} gives each casualty a MARCH state with what it rests on:
\begin{itemize}
  \item \textbf{M}assive bleeding: not controlled / controlled / tourniquet / control unknown.
  \item \textbf{A}irway: at risk if unresponsive, gurgling, snoring or gasping; open if talking.
  \item \textbf{R}espiration, \textbf{C}irculation (pale, cold, sweaty, thirsty, weak pulse = shock signs),
        \textbf{H}ead and hypothermia.
\end{itemize}
Responsiveness (AVPU) comes from what is said (``confused'', ``only to pain''), from a speaker who is the
casualty (``I'm hit''), or from vocalising heard on the clip. Distress comes from the casualty's own voice,
screams or groans heard, or the buddy's words. Each source can only raise it. The overall state is stable,
serious, critical or very critical (score 0--100).

\begin{table}[htbp]\centering\small
\caption{Threat windows (textbook orders of magnitude, to check with the medical mentor).}\label{tab:windows}
\begin{tabular}{lll}
\toprule
Threat & Window & Counted from\\
\midrule
Massive bleeding not controlled & 3--10 min & injury\\
Airway at risk & 3--5 min & now\\
Breathing problem & 5--30 min & now\\
Shock signs & 15--45 min & now\\
Bleeding, control unknown & 10--30 min & injury\\
Torso or junctional bleeding (no tourniquet possible) & 30--60 min & injury\\
Head injury or blast with reduced consciousness & 30--120 min & injury\\
Hypothermia (November, lying still) & 30--90 min & injury\\
Tourniquet on a limb & 120--360 min & tourniquet time\\
\bottomrule
\end{tabular}
\end{table}

The clock starts at the injury when it's said (``hit 6 minutes ago''), else before the tourniquet, else at the
first report. The read-back names the soonest threat when it's within 15 minutes (``Warning, critical within the
next 4 minutes: massive bleeding not controlled''). Tourniquet timers change colour at 60, 90 and 120 minutes.

\section{Mobility}

\code{mobility.py} gives each casualty one state: can walk, needs help, stretcher, trapped, or unknown.
\begin{itemize}
  \item \textbf{Said first}, in English or French: ``he can walk'', ``il ne peut pas marcher'', ``limping'', ``pinned
        under rubble''. The latest mention wins.
  \item \textbf{Otherwise inferred} and marked as such: unconscious, massive bleeding, breathing problem, a leg
        tourniquet, or a thigh or pelvis wound mean a stretcher. A lower-leg or foot wound means needs help.
        ``Conscious'' with no leg wound means can walk. Anything else is unknown, and a stretcher is planned.
  \item \textbf{Conflicts are flagged}: ``can walk'' said together with a leg tourniquet keeps what was said, shows
        the conflict, and plans the stretcher anyway.
\end{itemize}
Mobility decides between guidance to the CCP and a medic route (Chapter~\ref{ch:routing}).

\section{Body twin, impact graph, kit and team}

\paragraph{Posture} is estimated, never set by hand. Each clue adds points to some postures: what was said
weighs most; unconscious and breathing suggests the recovery position; a chest wound while conscious suggests
sitting up; a leg tourniquet suggests lying on the back. The scores become probabilities shown under a 3D body
(SMPL-X, built locally because its licence forbids redistribution). Injuries go onto the named body segment;
tourniquets are drawn as an orange band with their running time.

\paragraph{Impact graph.} \code{response.py} draws, from left to right: what was reported $\rightarrow$ body
region $\rightarrow$ structures possibly impacted (dashed, to check: thigh bleeding $\rightarrow$ femoral artery;
blast $\rightarrow$ eardrums, blast lung, concussion) $\rightarrow$ MARCH threat $\rightarrow$ what to bring. How it
happened points to where to look: a fall points to the neck, spine and pelvis; a penetrating wound to a possible
exit wound.

\paragraph{Body parts graph.} Body $\rightarrow$ 7 regions $\rightarrow$ 40 parts. Red means affected (said); orange
means no data yet, which is not the same as fine.

\paragraph{Kit and team.} Kit in MARCH order with quantities and who uses it: tourniquets for limb bleeding, packing
only for junctional wounds, chest seals for chest wounds, a pelvic binder for blast with leg wounds. The team is
the smallest that can carry: 4 bearers per stretcher (2 under 100 m, +2 relief beyond 500 m), escorts for walking
wounded, one medic per two urgent casualties, plus extraction tools for a trapped casualty. Every extra person is
one more target for drones.

\nanofig{casualty}{The casualty view: MARCH state with countdowns, the body twin, and the impact graph from what was said to
what to bring.}
{Medical decision-support infographic, 16:9, plain white background, clean flat vector, red, amber, green and grey status
colours. Left: a vertical MARCH strip with five rows ``M Massive bleeding: tourniquet'' (amber), ``A Airway: open''
(green), ``R Breathing: unknown'' (grey), ``C Circulation: shock signs'' (red), ``H Hypothermia risk'' (amber), each
with a small countdown bar; a label ``worse in 4–10 min''. Centre: a neutral grey 3D human body lying on its back,
seen from above at an angle, left thigh glowing red with an orange band labelled ``TQ 14:32 · 72 min''. Right: a
left-to-right graph of rounded boxes ``left leg bleeding'' → ``left thigh'' → ``femoral artery (check)'' (dashed) →
``M · massive bleeding'' → ``tourniquet, packing''. Small footer text ``decision support, not a diagnosis''. No gore.
Exactly these labels. Colour palette (use only these): navy \#0D1110, cyan \#1AA7C4, amber \#E09A2B, red \#E5484D, green \#3FA877, grey \#8B9A94, plus white for text on dark shapes.}
